% section: 形から決める第一手

\section{形から決める第一手}

ここでは,問題を見たときに最初に試す操作を,計算の流れとして整理する.
大切なのは,一つの方法に固執せず,計算が簡単になる方向へ進むことである.

\subsection{差なら和へ戻す}

\[
  a_{n+1}-a_n=f(n)
\]
なら,\,$n=1$ から $n-1$ まで足す.
\[
  a_n=a_1+\sum_{k=1}^{n-1}f(k).
\]
差分は,数列における積分に相当する.

\subsection{比なら積へ戻す}

\[
  \frac{a_{n+1}}{a_n}=g(n)
\]
なら,順に掛け合わせる.
\[
  a_n=a_1\prod_{k=1}^{n-1}g(k).
\]
分母が0になる可能性がある場合は,この変形を使える範囲を先に確認する.

\subsection{一次式なら定数を消す}

\[
  a_{n+1}=pa_n+q
\]
では,不動点 $\alpha$ を
\[
  \alpha=p\alpha+q
\]
で定める.\,$b_n=a_n-\alpha$ と置けば,
\[
  b_{n+1}=pb_n
\]
となる.定数を消すことで,一次漸化式を等比型へ変えている.

\subsection{既知の項が残るなら補助数列を作る}

\[
  a_{n+1}=pa_n+f(n)
\]
で $f(n)$ が邪魔をするなら,同じ種類の項を含む $g(n)$ を探し,
\[
  g(n+1)=pg(n)+f(n)
\]
を満たすようにする.\,$b_n=a_n-g(n)$ と置くと,
\[
  b_{n+1}=pb_n
\]
となる.これは,元の漸化式の特別な解を先に見つけて引く方法である.

\subsection{隣接する複数の項が線形に現れるなら}

\[
  a_{n+2}=pa_{n+1}+qa_n
\]
では,\,$a_n=\lambda^n$ を試す.
\[
  \lambda^2-p\lambda-q=0
\]
が得られ,この方程式の根から基本解を作る.重根なら $n\lambda^n$,
複素数根なら正弦・余弦の形が現れる.

\subsection{分数の中に数列があるなら}

\[
  a_{n+1}=\frac{pa_n+q}{ra_n+s}
\]
では,まず不動点を探す.逆数だけで一次式になる場合は
$b_n=1/a_n$ を使う.不動点が二つある場合は,
\[
  b_n=\frac{a_n-\alpha}{a_n-\beta}
\]
を試す.この比が等比型になれば,分数の反復が積へ変わる.

\subsection{どの形にも入らないなら}

そんなときは特殊関数型漸化式であるか,突飛な発想を求められるものである.
一旦ご飯でも食べて改めて取り組もう.
ご飯の時間でないのなら,今までの数学の知識を総動員して似たものを探そう.